\pg{358} O ente necessario, infinito, semplicissimo, é só acto, sem nenhuma sombra de potencialidade. Ora o que é só acto, sem sombra de potencialidade, é immutavel. Logo o ente \emph{necessario}, \emph{infinito}, \emph{simplicissimo}, \emph{é immutavel} \footnote{O ente \emph{immutavel} --- é \emph{necessario}, porque, se fosse \emph{contingente}, receberia de outro ente a existencia, e pelo facto de passar da potencia para o acto, já não seria \emph{immutavel}. --- é \emph{infinito}, porque, se fosse \emph{finito}, estaria em potencia para uma perfeição ulterior, e deixaria de ser \emph{immutavel}. --- é \emph{simplicissimo}, porque, se fosse \emph{composto}, tambem deixaria de ser \emph{immutavel}. --- O ente \emph{necessario}, \emph{infinito}, \emph{simplicissimo}, é \emph{immutavel}; porque, --- se o ente \emph{necessario} fosse \emph{mutavel}, ou adquiriria uma perfeição que não tinha, ou perderia uma perfeição que tinha, e assim não seria \emph{necessario}, mas \emph{contingente} em relação á perfeição adquirida ou perdida, --- se o ente \emph{infinito} fosse \emph{mutavel}, tambem ou perderia uma perfeição, e não seria \emph{infinito}, mas \emph{finito}, --- se o ente \emph{simplicissimo} fosse \emph{mutavel}, passaria da potencia para o acto, e já não seria \emph{simplicissimo}, mas \emph{composto} (Cf. \emph{Sum}. \emph{Th}., l. c.). Como se vê, \emph{acto puro}, \emph{ente necessario}, \emph{ente infinito}, \emph{ente simplicissimo}, \emph{ente immutavel}, são attributos que exprimem uma e a mesma coisa, e por isso a um ente, ao qual convenha um d'esses attributos, convêm todos, e de um d'elles pode legitimamente deduzir-se outro.}. §8. \textbf{O ente mutavel é contingente, finito, composto; o ente contingente, finito, composto, é mutavel.}

a) \emph{O ente mutavel é contingente, finito, composto.} --- Ente mutavel é o que está sujeito a mudanças, e por isso não é acto puro. Ora o que não é acto puro é contingente, finito, composto. Logo o ente \emph{mutavel é contingente, finito, composto}.

b) \emph{O ente contingente, finito, composto, é mutavel.} --- Todo o ente contingente, finito, composto, não é acto puro, mas tem em si alguma potencialidade. Ora o que tem potencialidade é mutavel. Logo o ente \emph{contingente, finito, composto, é mutavel} \footnote{O ente \emph{mutavel} --- é \emph{contingente}, porque, na mudança, adquire ou perde uma perfeição, e assim é \emph{contingente} em relação á perfeição adquirida ou perdida, --- é \emph{finito}, porque, se, na mudança, adquire ou perde uma perfeição, não possue todas as perfeições, e por isso é \emph{finito}, --- é \emph{composto}, porque, quando \emph{muda}, passa da potencia para o acto, e por isso consta d'esses dois elementos. --- O ente \emph{contingente, finito, composto, é mutavel}. É mutavel o ente \emph{contingente}; porque, sendo \emph{contingente}, passou e pode passar da potencia para o acto, e assim mudar-se. É mutavel o ente \emph{finito}; porque, sendo \emph{finito}, pode adquirir uma perfeição ulterior, e assim passar da potencia para o acto. É mutavel o ente \emph{composto}; porque, sendo composto de potencia e de acto, pode passar de um estado para outro. (\emph{Sum}. \emph{Th}., l. c., a. 2).}.

§9. \textbf{O ente mutavel suppõe o ente immutavel.} --- Na ordem \emph{ontologica}, ou da \emph{actualidade}, o ente \emph{actual} deve preceder o ente \emph{possivel}; porque este só pode adquirir o acto pela acção de um